% ============================================
% C. Natural Sample-Derived Communities
% ============================================
\subsection*{C. Natural Sample-Derived Communities}
\label{sec:suppl:natural_communities}

To test whether the patterns of community-level selection (CLS) observed in synthetic communities generalize to more ecologically relevant settings, we performed coalescence experiments using communities derived from natural environmental samples.

\subsubsection*{Experimental approach}

We collected 6 natural samples from diverse environmental sources (soil, compost, and decomposing organic matter). From each sample, we established bacterial communities through serial growth-dilution cycles in laboratory media (Supplementary Fig.~S28A):

\begin{enumerate}
    \item \textbf{Natural sample collection}: Six environmental samples ($N_{\text{Samples}} = 6$) were collected from distinct microhabitats.
    \item \textbf{Community assembly}: Each sample was inoculated into growth media and subjected to 7 rounds of serial growth-dilution (24h cycles), allowing communities to stabilize under laboratory conditions ($N_{\text{Communities}} = 6$).
    \item \textbf{Coalescence}: Stabilized communities were mixed pairwise at 50:50 ratio and allowed to stabilize through 7 additional rounds of serial growth-dilution ($N_{\text{Coalesced}} = 16$).
    \item \textbf{Measurement}: Community composition was characterized via 16S rRNA amplicon sequencing, along with optical density (OD) and pH measurements.
\end{enumerate}

This approach yields communities with natural evolutionary histories and ecological interactions, in contrast to our synthetic communities assembled from isolated strains.

\subsubsection*{Results}

We analyzed coalescence outcomes using the same vector decomposition framework applied to synthetic communities. The results reveal striking similarities between natural and synthetic systems (Supplementary Fig.~S28B,C):

\textbf{Phase diagram analysis}: Natural sample-derived communities exhibited the same qualitative patterns as synthetic communities across all three nutrient conditions (Nutr$-$, Base, Nutr+). Points cluster along the axes (indicating CLS outcomes) rather than along the diagonal (which would indicate mixing), consistent with one-sided selection dynamics.

\textbf{Outcome distributions}: The fraction of CLS outcomes increased with nutrient availability, mirroring the trend observed in synthetic communities:
\begin{itemize}
    \item In Nutr$-$ medium, outcomes were distributed across CLS, Mixing, and Restructuring categories
    \item In Base and Nutr+ media, CLS dominated the outcome distribution
\end{itemize}

These results demonstrate that community-level selection is a robust phenomenon that emerges across different community assembly histories---from precisely controlled synthetic consortia to complex, naturally-evolved microbial assemblages.

\subsubsection*{Interpretation}

The consistency between natural and synthetic community experiments suggests that:
\begin{enumerate}
    \item CLS does not depend on specific strain identities or artificially simplified community structures
    \item The underlying mechanism---correlated species retention within assembled communities---operates in ecologically realistic settings
    \item Nutrient availability modulates interaction strength and thus CLS frequency, regardless of community origin
\end{enumerate}

% Figure moved to supplementary_sections/figures.tex (Supplementary Fig.~S28)
